\section{Scope, repository audit, and proof boundary}

\subsection{The checked core and the research layer}

At the inspected snapshot
\[
  \code{b88e843c31b50dfe461bfbd0380166ca422cd3ab},
\]
the directory \code{Logic/PrimitiveRecursiveArithmetic} has two deliberately
different layers.  The first is \code{ConciseIPRA}, a Lean development of
small systems bi-interpretable with ordinary intuitionistic primitive
recursive arithmetic.  The second is a catalog of fifty-eight delivered
research packets.  The repository warns that these packets were filed as
delivered and were not collectively rebuilt or verified.  In particular, no
packet currently supplies a completed derivation of its main mathematical
theorem in the object type
\[
       \Gamma\;\vdash[\PRAAx]\;\varphi.
\]
This report takes that distinction as a design constraint, not as a clerical
detail.

The object language actually available in the repository is unusually useful
for this project.  The inductive type \code{PR k} contains zero, successor,
projections, composition, and primitive recursion.  Terms \code{PTm n} are
applications of those symbols to variables.  The axiom family \PRAAx\ contains
the defining equations, successor separation and injectivity, and induction
for quantifier-free formulas.  The theorem \code{PRA.qfInd} is therefore the
principal internal engine.  The theorem \code{pra\_iff\_native} transports an
already completed \PRAAx\ derivation to the tiny native calculus; it is not a
device for turning an arbitrary semantic Lean proof into such a derivation.

\subsection{What ``embedding mathematics in iPRA'' means}

A mathematically serious claim of an iPRA proof should contain four separate
objects.

\begin{enumerate}[label=\textbf{(\arabic*)},leftmargin=2.2em]
\item A finite numerical presentation of the input, including every modulus,
      separation constant, height cutoff, chart, or cover used by the proof.
\item Fixed primitive-recursive algorithms for the finite computations.
      ``Fixed'' matters: a universal evaluator taking an arbitrary PR program
      code as data is not itself primitive recursive.
\item A quantifier-free invariant and an object-language derivation from
      \PRAAx.  A native theorem about Lean's \code{Nat}, even a constructive
      one, is only a semantic lemma until this derivation is built.
\item An external interpretation theorem connecting the numerical
      presentation to the familiar mathematical statement.
\end{enumerate}

The order is important.  One proves the arithmetic certificate first and
transports it through the repository's bi-interpretation only afterward.
This prevents three common category errors: treating type checking of a goal
as a proof, treating a finite test run as a universal theorem, and treating a
classical existence theorem as if it supplied a primitive-recursive witness.

\subsection{Recurring difficulties}

The proof packets reveal a stable list of obstacles.

\begin{longtable}{p{.20\linewidth}p{.34\linewidth}p{.36\linewidth}}
\toprule
\textbf{Obstacle} & \textbf{Typical invalid move} & \textbf{Required repair}\\
\midrule
\endhead
Unbounded precision &
Induct on a proposition saying that an identity holds at every future
precision. &
Fix a free requested precision, compute the finite observations it demands,
prove a bounded invariant, then universally introduce the request.\\

Discontinuous exact choice &
Select the sign, root, eigenvector, rotation axis, or Puiseux branch of
arbitrary Cauchy data. &
Use overlapping branches, finite relational outputs, projective data, or
supply a positive separation certificate.\\

Abstract compactness &
Invoke compactness or completeness without exhibiting a finite cover. &
Provide a primitive-recursive net, a Lebesgue-number-type margin, and a
finite refinement verifier.\\

Finiteness without a cutoff &
Enumerate all rational or integral points from a theorem saying only that
there are finitely many. &
Separate the bounded verifier from the missing height or covering
certificate producer.\\

Ineffective constants &
Use an asymptotic theorem whose constants are not given effectively. &
Make the constants part of the input certificate and prove a verifier theorem
with the correct quantifier order.\\

Semantic/object mismatch &
Prove the intended theorem in ambient Lean and call it an iPRA proof. &
Compile the algorithm to \code{PR}, express its invariant as an
\code{Fm PTm}, and construct a \code{Prf PRAAx} term.\\
\bottomrule
\end{longtable}

These obstacles are not interchangeable.  A missing analytic estimate is not
fixed by better coding; a discontinuous selector is not fixed by a faster
search; and a semantic Lean proof is not made internal by adding more test
cases.

\subsection{Contribution of this report}

The current repository already contains several compatible frameworks:
certificate-and-modulus interfaces, finite observations, evidence-indexed
outputs, uniform certificates, and the newly placed Certified Finitary
Presentation report.  The present contribution is an integration layer.

\begin{itemize}[leftmargin=1.7em]
\item A \emph{contract-DAG} packages a whole proof as a fixed finite acyclic
      network of local certified transducers.
\item A reverse pass computes all precision and observation demands.
\item A forward bounded trace evaluates the nodes.
\item One quantifier-free trace invariant composes the local proof objects.
\item Relational outputs are handled by a finite-horizon compatibility
      theorem rather than by hidden choice.
\item A packet-conversion protocol identifies, for every theorem, the exact
      missing producer and the exact reusable verifier.
\end{itemize}

The result is deliberately modular.  It does not make the general Faltings or
Siegel theorem effective, does not manufacture the missing constants in
ineffective prime-distribution estimates, and does not produce exact Puiseux
branches from arbitrary Cauchy coefficient names.  It does show precisely
what can be proved once the appropriate finite certificate is supplied and
how that proof should be assembled in the repository's actual object theory.
